\documentclass[spanish,letterpaper,oneside]{article}

\def\documenttitle {S3. Revolucion francesa, independencia mexicana y
Sentimientos de la Nacion}
\def\documentsubtitle {Revolucion francesa e independencia mexicana}
\def\documentsubject {S3 - Teoria del Estado y Constitucion}
\def\documentauthor {Martin Jonathan de la Cruz}
\def\coursename {Teoria del Estado y Constitucion}
\def\coursecode {020}
\def\universityname {Universidad Abierta y a Distancia de Mexico}
\def\universityfaculty {Licenciatura en Derecho}
\def\universitydepartment {Teoria del Estado y Constitucion}
\def\universitydepartmentimage {departamentos/UnADM}
\def\universitydepartmentimagecfg {height=1.57cm}
\def\universitylocation {Roma Norte, Ciudad de Mexico}
\def\authortable {
  \begin{tabular}{ll}
    \textbf{Alumno:} & Martin Jonathan de la Cruz \\
    Matricula: & ES2611202040 \\
    Figura docente: & Ana Lilian Sanchez Chavez \\
    Semestre/Bloque: & 2 / 2 \\
    Estado: & Borrador para revision \\
    \multicolumn{2}{l}{Fecha de realizacion: \today} \\
    \multicolumn{2}{l}{\universitylocation}
  \end{tabular}
}

\input{template}
\usepackage{helvet}
\renewcommand{\familydefault}{\sfdefault}
\usepackage{array}
\usepackage{booktabs}
\usepackage{longtable}
\usepackage{calc}
\usepackage{tikz}
\usetikzlibrary{arrows.meta,positioning,calc}
\setcitestyle{authoryear,open={(},close={)}}
\providecommand{\tightlist}{\setlength{\itemsep}{0pt}\setlength{\parskip}{0pt}}
\providecommand{\pandocbounded}[1]{#1}
\providecommand{\passthrough}[1]{#1}
\setlength{\emergencystretch}{3em}

\def\coverwatermarkenabled {true}
\def\coverwatermarkimage {img/departamentos/UnADM.pdf}
\def\coverwatermarkopacity {0.16}
\def\coverwatermarkwidth {0.72\paperwidth}
\def\coverwatermarkxshift {0cm}
\def\coverwatermarkyshift {-0.35cm}
\newcommand{\insertcoverwatermark}{%
  \ifthenelse{\equal{\coverwatermarkenabled}{true}}{%
    \AddToShipoutPictureBG*{%
      \begin{tikzpicture}[remember picture,overlay]
        \node[opacity=\coverwatermarkopacity,inner sep=0pt,
          xshift=\coverwatermarkxshift,yshift=\coverwatermarkyshift]
          at (current page.center) {%
            \includegraphics[width=\coverwatermarkwidth]{\coverwatermarkimage}%
          };
      \end{tikzpicture}%
    }%
  }{}%
}

\begin{document}
\insertcoverwatermark
\templatePortrait
\templatePagecfg
\onehalfspacing
\templateIndex
\templateFinalcfg

\textbf{UnADM \textbar{} Licenciatura en Derecho \textbar{} Teoria del
Estado y Constitucion \textbar{} Grupo 020}

\textbf{Borrador anticipado para revision.} Entrega docente: 18 octubre
de 2026, 23:55 HC. Falta identidad, formato y cotejo de ediciones
primarias. Asistencia de GitHub Copilot en redaccion y organizacion; no
enviado. Las relaciones propuestas son afinidades conceptuales, no
prueba de influencia directa documentada.

\section{Problema de analisis}\label{problema-de-analisis}

La Revolucion francesa y la independencia mexicana cuestionaron ordenes
politicos diferentes. La primera se desarrollo en la crisis del Antiguo
Regimen y redefinio legitimidad, igualdad juridica y soberania nacional.
La segunda incluyo la ruptura de una relacion colonial y proyectos
diversos sobre organizacion, religion y justicia social. Compararlas
exige reconocer ideas compartidas sin reducir el proceso mexicano a una
copia de Francia.

\section{Rousseau: soberania y voluntad
general}\label{rousseau-soberania-y-voluntad-general}

En \emph{Del contrato social}, Rousseau situa la soberania en la
voluntad general y rechaza que pueda alienarse o dividirse como
titularidad soberana (libro II, caps. 1-2). No equivale a que todas las
funciones deban concentrarse en una persona. El punto 5 de
\emph{Sentimientos de la Nacion}, presentado por Morelos el 14 de
septiembre de 1813, coloca la soberania en el pueblo y organiza su
deposito mediante representacion. La afinidad consiste en desplazar el
fundamento del poder de la voluntad de un monarca a la comunidad
politica.

No son construcciones identicas: Rousseau cuestiona que la soberania
misma pueda representarse (\emph{Del contrato social}, libro III, cap.
15), mientras Morelos formula instituciones representativas. La
comparacion permite distinguir titularidad popular y ejercicio
institucional, sin afirmar identidad entre ambos modelos ni probar que
Morelos adoptara directamente cada tesis de Rousseau.

\section{Montesquieu: limites y division
funcional}\label{montesquieu-limites-y-division-funcional}

En \emph{Del espiritu de las leyes} (libro XI, cap. 6), Montesquieu
argumenta que la libertad politica requiere evitar que funciones
legislativa, ejecutiva y de juzgar se acumulen de forma que permita
arbitrariedad. El punto 6 de \emph{Sentimientos} propone dividir los
poderes legislativo, ejecutivo y judicial en cuerpos compatibles para
ejercerlos. Existe afinidad en la necesidad de distribuir funciones y
limitar concentracion, no una demostracion documental de recepcion
textual.

La division funcional y la soberania popular responden a preguntas
distintas: quien es titular de la autoridad y como se organiza su
ejercicio. Confundirlas produce la falsa contradiccion de pensar que
dividir poderes necesariamente destruye la unidad de la comunidad
politica.

\section{Revolucion francesa y proyecto
mexicano}\label{revolucion-francesa-y-proyecto-mexicano}

La Declaracion francesa de 1789 expresa igualdad juridica (art. 1),
soberania nacional (art. 3), ley vinculada a voluntad general (art. 6) y
derechos garantizados con separacion de poderes (art. 16). En Morelos,
los puntos 5-6 articulan soberania y organizacion; el 15 dispone
proscribir esclavitud y distinciones de castas; el 12 incorpora
moderacion de opulencia e indigencia y mejora de condiciones del
trabajador. Estos elementos permiten leer una exigencia de justicia
social que no se agota en reproducir la Declaracion francesa.

Tambien hay diferencias: el punto 2 de \emph{Sentimientos} afirma
exclusividad de la religion catolica, tension con concepciones
contemporaneas de pluralismo. El punto 1 declara independencia de Espana
y de toda otra nacion, objetivo propio de emancipacion colonial. Por
ello no resulta riguroso presentar el proyecto como constitucionalismo
liberal homogeneo ni juzgarlo sin contexto historico.

\section{Conclusion critica
propuesta}\label{conclusion-critica-propuesta}

Las ideas europeas ofrecen categorias para pensar soberania, libertad y
controles, pero su incorporacion se transforma ante condiciones
mexicanas de colonialidad, castas y desigualdad. La aportacion de
\emph{Sentimientos de la Nacion} consiste en vincular legitimidad
popular y organizacion del poder con emancipacion y demandas sociales.
La relevancia actual no radica en repetir formulas: requiere evaluar si
las instituciones controlan arbitrariedad y permiten ejercer derechos.
Esta lectura reconoce afinidades y diferencias; demostrar influencia
historica directa exigiria evidencia de circulacion y recepcion que este
borrador no aporta.

\section{Referencias primarias para cotejo de
edicion}\label{referencias-primarias-para-cotejo-de-edicion}

Asamblea Nacional Constituyente de Francia. (1789, 26 de agosto).
\emph{Declaracion de los Derechos del Hombre y del Ciudadano}.

Montesquieu, C. L. de Secondat. (1748). \emph{Del espiritu de las leyes}
(libro XI, cap. 6). Cotejar traduccion, editor y ano de la edicion
efectivamente utilizada.

Morelos y Pavon, J. M. (1813, 14 de septiembre). \emph{Sentimientos de
la Nacion} (puntos 1, 2, 5, 6, 12 y 15). Cotejar reproduccion
documental.

Rousseau, J.-J. (1762). \emph{Del contrato social} (libro II, caps. 1-2;
libro III, cap. 15). Cotejar traduccion y edicion consultada.

Sanchez Chavez, A. L. (2026). \emph{Planificacion de actividades:
Semanas 2 y 3, Teoria del Estado y Constitucion} {[}Consigna de aula{]}.
Universidad Abierta y a Distancia de Mexico.

\textbf{Nota:} fechas originales no sustituyen datos APA de la edicion
consultada. No incluir referencias de obras no leidas como si hubieran
sido verificadas. Completar la bibliografia al hacer el cotejo final.

\end{document}